\subsection{Geometry-Input Ablation in Closed Loop}
The similar K12 discovery counts motivate a separate execution comparison using the two already frozen networks. The 512 unaffected baseline episodes reproduce all outcomes, 28,593 recorded path states and 28,081 planner decisions exactly. The comparison therefore isolates which trained network supplies the priority for Learned and Hybrid.

Table~\ref{tab:geometryablation} shows that removing the clearance channels increases Learned completion from 99 to 117 of 128 tasks and reduces margin incursions from nine to two. The paired comparison contains 23 No-geometry-only completions and five Geometry-only completions. Geometry minus No-geometry has a completion-rate difference of $-14.06$ percentage points [$-21.09$,$-7.03$]. The largest difference occurs under crosswind, where No-geometry completes 26 of 32 tasks compared with 12 for Geometry. For Hybrid, the corresponding overall counts are 92 and 89; the difference is $-2.34$ percentage points [$-12.50$,9.38].

\begin{table}[!t]
\centering\small
\caption{Frozen-model geometry-input ablation. Each row contains the same 128 map--condition tasks. Margin denotes entry into inflated obstacle geometry.}
\label{tab:geometryablation}
\begin{tabular}{llrr}
\toprule
Schedule & Input & Goals/128 & Margin \\
\midrule
Learned & Geometry & 99 & 9 \\
Learned & No geometry & 117 & 2 \\
Hybrid & Geometry & 89 & 10 \\
Hybrid & No geometry & 92 & 9 \\
\bottomrule
\end{tabular}
\end{table}

The geometry model predicts the nominal finite-bank label more accurately and tends to select a lower-cost verified primary, yet this advantage does not transfer to mission completion. The explicit validator already supplies geometric acceptance, while the learned dynamic-state, goal and cost features can favor different maneuvers among candidates that pass. This result locates the strongest demonstrated benefit in learning continuation-query priorities; it does not establish an execution benefit from the additional clearance-grid representation. The ablation uses the original test maps and frozen weights, so it adds a representation comparison rather than a new generalization split.
